\documentclass[11pt]{article}
\usepackage[a4paper,margin=24mm]{geometry}
\usepackage{fontspec}
\usepackage{xcolor}
\usepackage{amsmath,amssymb}
\usepackage{booktabs,longtable,array}
\usepackage[hidelinks]{hyperref}
\setmainfont[Path=fonts/,Extension=.otf,UprightFont=*-Regular,BoldFont=*-Bold]{NotoSerifCJKkr}
\definecolor{reviewercolor}{RGB}{38,82,126}
\definecolor{responsecolor}{RGB}{36,113,78}
\definecolor{revisioncolor}{RGB}{145,82,24}
\setlength{\parindent}{0pt}
\setlength{\parskip}{0.55em}
\setlength{\fboxsep}{7pt}
\setlength{\emergencystretch}{3em}
\renewcommand{\familydefault}{\rmdefault}
\raggedright
\newcommand{\fieldheading}[2]{\par\medskip\noindent{\color{#1}\rule{3pt}{1.25em}}\hspace{0.55em}\textbf{#2}\par\smallskip}
\newenvironment{reviewitem}[1]{\subsection*{#1}}{\par\medskip\hrule\bigskip}
\newenvironment{reviewercomment}{\fieldheading{reviewercolor}{1. 리뷰어 코멘트}\begin{quote}}{\end{quote}}
\newenvironment{authorresponse}{\fieldheading{responsecolor}{2. 저자 답변}}{\par}
\newenvironment{manuscriptrevision}[1]{\fieldheading{revisioncolor}{3. 논문 반영 결과}\textbf{수정 위치:} #1\par\smallskip}{\par}
\newcommand{\manuscripttitle}{연속 인과관계(CoC) 주석의 시간·조건 일관성 검사 프레임워크: 적용 가능성과 근거 추적}
\newcommand{\manuscriptid}{확인 필요 (제공 기록에서 확인되지 않음)}
\newcommand{\revisionround}{확인 필요 (개정 소스 버전: revision-v01)}
\newcommand{\responsedate}{2026년 10월 12일}
% Generated from current manuscript auxiliary labels.
\expandafter\def\csname mloc@sec:intro\endcsname{1절(1쪽)}
\expandafter\def\csname mloc@fig:overview-pipeline\endcsname{그림 1(2쪽)}
\expandafter\def\csname mloc@sec:related\endcsname{2절(2쪽)}
\expandafter\def\csname mloc@sec:related-reasoning\endcsname{2.1절(2쪽)}
\expandafter\def\csname mloc@sec:related-rules\endcsname{2.2절(2쪽)}
\expandafter\def\csname mloc@sec:related-nl\endcsname{2.3절(2쪽)}
\expandafter\def\csname mloc@sec:related-rv\endcsname{2.4절(3쪽)}
\expandafter\def\csname mloc@sec:related-position\endcsname{2.5절(3쪽)}
\expandafter\def\csname mloc@sec:data\endcsname{3절(3쪽)}
\expandafter\def\csname mloc@sec:data-overview\endcsname{3.1절(3쪽)}
\expandafter\def\csname mloc@sec:provenance\endcsname{3.2절(3쪽)}
\expandafter\def\csname mloc@sec:review-protocol\endcsname{3.3절(3쪽)}
\expandafter\def\csname mloc@sec:problem\endcsname{4절(3쪽)}
\expandafter\def\csname mloc@tab:related-position\endcsname{표 1(4쪽)}
\expandafter\def\csname mloc@tab:evaluation-data\endcsname{표 2(4쪽)}
\expandafter\def\csname mloc@tab:provenance\endcsname{표 3(4쪽)}
\expandafter\def\csname mloc@sec:notation\endcsname{4.1절(4쪽)}
\expandafter\def\csname mloc@sec:error-types\endcsname{4.2절(4쪽)}
\expandafter\def\csname mloc@eq:hold-conflict\endcsname{식 2(4쪽)}
\expandafter\def\csname mloc@eq:missing-evidence\endcsname{식 6(4쪽)}
\expandafter\def\csname mloc@tab:model-scope\endcsname{표 4(5쪽)}
\expandafter\def\csname mloc@sec:semantics\endcsname{4.3절(5쪽)}
\expandafter\def\csname mloc@eq:active-update\endcsname{식 9(5쪽)}
\expandafter\def\csname mloc@sec:boundary-semantics\endcsname{4.4절(5쪽)}
\expandafter\def\csname mloc@tab:boundary-policy\endcsname{표 5(5쪽)}
\expandafter\def\csname mloc@sec:issue-definition\endcsname{4.5절(5쪽)}
\expandafter\def\csname mloc@tab:parser-rules\endcsname{표 6(6쪽)}
\expandafter\def\csname mloc@eq:response-contract\endcsname{식 13(6쪽)}
\expandafter\def\csname mloc@sec:method\endcsname{5절(6쪽)}
\expandafter\def\csname mloc@sec:workflow\endcsname{5.1절(6쪽)}
\expandafter\def\csname mloc@sec:extraction\endcsname{5.2절(6쪽)}
\expandafter\def\csname mloc@sec:models\endcsname{5.3절(6쪽)}
\expandafter\def\csname mloc@sec:queries\endcsname{5.3절(6쪽)}
\expandafter\def\csname mloc@sec:trace\endcsname{5.3절(6쪽)}
\expandafter\def\csname mloc@lst:condition-step\endcsname{목록 1(6쪽)}
\expandafter\def\csname mloc@fig:method-pipeline\endcsname{그림 2(7쪽)}
\expandafter\def\csname mloc@tab:worked-example\endcsname{표 7(7쪽)}
\expandafter\def\csname mloc@tab:model-variables\endcsname{표 8(7쪽)}
\expandafter\def\csname mloc@fig:condition-observer\endcsname{그림 3(7쪽)}
\expandafter\def\csname mloc@lst:queries\endcsname{목록 2(7쪽)}
\expandafter\def\csname mloc@tab:query-meaning\endcsname{표 9(7쪽)}
\expandafter\def\csname mloc@sec:boundary-implementation\endcsname{5.4절(7쪽)}
\expandafter\def\csname mloc@sec:model-scope\endcsname{5.5절(8쪽)}
\expandafter\def\csname mloc@sec:experiments\endcsname{6절(8쪽)}
\expandafter\def\csname mloc@sec:legacy-test\endcsname{6.1절(8쪽)}
\expandafter\def\csname mloc@sec:scenario-design\endcsname{6.2절(8쪽)}
\expandafter\def\csname mloc@sec:retest-design\endcsname{6.3절(8쪽)}
\expandafter\def\csname mloc@sec:boundary-design\endcsname{6.4절(8쪽)}
\expandafter\def\csname mloc@sec:bounded-design\endcsname{6.5절(8쪽)}
\expandafter\def\csname mloc@tab:comparison-paths\endcsname{표 10(9쪽)}
\expandafter\def\csname mloc@tab:scenario-variants\endcsname{표 11(9쪽)}
\expandafter\def\csname mloc@sec:human-design\endcsname{6.6절(9쪽)}
\expandafter\def\csname mloc@sec:aux-design\endcsname{6.7절(9쪽)}
\expandafter\def\csname mloc@tab:regression\endcsname{표 12(9쪽)}
\expandafter\def\csname mloc@sec:results\endcsname{7절(9쪽)}
\expandafter\def\csname mloc@sec:regression-results\endcsname{7.1절(9쪽)}
\expandafter\def\csname mloc@sec:first-results\endcsname{7.2절(9쪽)}
\expandafter\def\csname mloc@tab:first-results\endcsname{표 13(10쪽)}
\expandafter\def\csname mloc@tab:failures\endcsname{표 14(10쪽)}
\expandafter\def\csname mloc@sec:retest-results\endcsname{7.3절(10쪽)}
\expandafter\def\csname mloc@sec:agreement-results\endcsname{7.4절(10쪽)}
\expandafter\def\csname mloc@tab:retest\endcsname{표 15(10쪽)}
\expandafter\def\csname mloc@tab:implementation\endcsname{표 16(10쪽)}
\expandafter\def\csname mloc@tab:boundary-results\endcsname{표 17(10쪽)}
\expandafter\def\csname mloc@tab:bounded-results\endcsname{표 18(10쪽)}
\expandafter\def\csname mloc@sec:boundary-results\endcsname{7.5절(10쪽)}
\expandafter\def\csname mloc@sec:bounded-results\endcsname{7.6절(10쪽)}
\expandafter\def\csname mloc@tab:agreement\endcsname{표 19(11쪽)}
\expandafter\def\csname mloc@tab:reviewer-distribution\endcsname{표 20(11쪽)}
\expandafter\def\csname mloc@sec:human-results\endcsname{7.7절(11쪽)}
\expandafter\def\csname mloc@sec:trajectory-results\endcsname{7.8절(11쪽)}
\expandafter\def\csname mloc@sec:discussion\endcsname{8절(11쪽)}
\expandafter\def\csname mloc@sec:discussion-results\endcsname{8.1절(11쪽)}
\expandafter\def\csname mloc@sec:review-use\endcsname{8.2절(11쪽)}
\expandafter\def\csname mloc@tab:disagreement-examples\endcsname{표 21(12쪽)}
\expandafter\def\csname mloc@tab:trajectory-cases\endcsname{표 22(12쪽)}
\expandafter\def\csname mloc@tab:review-actions\endcsname{표 23(12쪽)}
\expandafter\def\csname mloc@sec:threats\endcsname{8.3절(12쪽)}
\expandafter\def\csname mloc@sec:future-work\endcsname{8.4절(12쪽)}
\expandafter\def\csname mloc@sec:conclusion\endcsname{9절(12쪽)}
\expandafter\def\csname mloc@sec:appendix-trajectory\endcsname{9절(12쪽)}
\expandafter\def\csname mloc@tab:limitations\endcsname{표 24(13쪽)}
\expandafter\def\csname mloc@tab:trajectory-details\endcsname{표 25(13쪽)}

% Generated from exact excerpts in the current main.pdf.
\expandafter\def\csname quotepages@R1-m3\endcsname{12쪽}
\expandafter\def\csname quotepages@R1-m2\endcsname{3쪽}
\expandafter\def\csname quotepages@R1-m1\endcsname{3쪽}
\expandafter\def\csname quotepages@R1-M6\endcsname{3쪽}
\expandafter\def\csname quotepages@R1-M5\endcsname{11쪽}
\expandafter\def\csname quotepages@R1-M4\endcsname{10쪽}
\expandafter\def\csname quotepages@R1-M3\endcsname{11쪽}
\expandafter\def\csname quotepages@R1-M2\endcsname{6쪽}
\expandafter\def\csname quotepages@R1-M1\endcsname{12쪽}
\expandafter\def\csname quotepages@R2-7\endcsname{2쪽}
\expandafter\def\csname quotepages@R2-6\endcsname{3쪽}
\expandafter\def\csname quotepages@R2-5\endcsname{1쪽}
\expandafter\def\csname quotepages@R2-4\endcsname{2쪽}
\expandafter\def\csname quotepages@R2-3\endcsname{1쪽}
\expandafter\def\csname quotepages@R2-2\endcsname{1쪽}
\expandafter\def\csname quotepages@R2-1.2\endcsname{1쪽}
\expandafter\def\csname quotepages@R2-1.1\endcsname{1쪽}
\expandafter\def\csname quotepages@R3-11\endcsname{11쪽}
\expandafter\def\csname quotepages@R3-10\endcsname{9쪽}
\expandafter\def\csname quotepages@R3-9\endcsname{11쪽}
\expandafter\def\csname quotepages@R3-8\endcsname{10쪽}
\expandafter\def\csname quotepages@R3-7\endcsname{10쪽}
\expandafter\def\csname quotepages@R3-6\endcsname{6쪽}
\expandafter\def\csname quotepages@R3-5\endcsname{6쪽}
\expandafter\def\csname quotepages@R3-4\endcsname{5쪽}
\expandafter\def\csname quotepages@R3-3\endcsname{4쪽}
\expandafter\def\csname quotepages@R3-2\endcsname{3쪽}
\expandafter\def\csname quotepages@R3-1\endcsname{3쪽}

\newcommand{\mloc}[1]{\csname mloc@#1\endcsname}
\newcommand{\quotepages}[1]{\csname quotepages@#1\endcsname}
\newcommand{\responsefront}[2]{%
\begin{center}{\LARGE\bfseries 심사 의견 답변서\par}\vspace{1em}{\large\bfseries\manuscripttitle\par}\end{center}
\vspace{0.8em}
\begin{tabular}{@{}ll@{}}논문 접수번호: & \manuscriptid \\ 수정 차수: & \revisionround \\ 답변일: & \responsedate\end{tabular}
\section*{답변에 앞서}
편집위원장님과 심사위원님께,

연속 CoC의 검사 범위, 의미 변환과 평가 근거를 구체화할 수 있도록 의견을 주셔서 감사합니다.
이번 개정은 자연 오류 성능 대신 시간·조건 관계의 검사 적용 가능성과 근거 추적으로 주장을 제한했습니다.
기존 통제 시나리오의 최초 결과와 개발 후 동일 사례 재시험을 구분하고, 81건의 장면 경계 시험과 기존 사람 라벨의 AC1 재분석을 본문에 추가했습니다.
선별 없는 유한 직적 전수 시험은 고정 XML의 step 코드와 별도 수식 판정기를 대조한 구현 적합성 근거입니다.
두 경로에 같은 작성자가 관여했으므로 외부 독립 시험이나 자연 의미 정답 검증을 완료했다고 주장하지 않습니다. 사람 재평가도 수행하지 않았습니다.
개정 원고의 파란색은 원 저장소의 최초 제출본 대비 이번에 추가·수정한 내용을 나타냅니다.
검정색 clean 원고는 동일 내용입니다. 아래 절·쪽·표 번호는 이번 개정 원고를 기준으로 합니다.
접수번호와 실제 심사 수정 차수는 기록이 없어 확인 필요로 남겼습니다.

\section*{심사위원 #1}
\textbf{총평(요지)}\par
\begin{quote}#2\end{quote}}


\begin{document}

\responsefront{2}{서론의 문장과 구조, 전반적인 국문 표현, UPPAAL 명칭, 연구 질문과 기여의 관계 및 관련 연구 대비 기술적 차이를 명확히 할 필요가 있다. 원문에서 반복된 의견은 한 번만 수록하고 문장 수정 1.1·1.2는 각각 답변한다.}

\clearpage

\begin{reviewitem}{R2-1.1. CoC 간 양립 문장}

\begin{reviewercomment}

기존 문장: ``문제는 같은 장면의 CoC들이 개별적으로는 자연스러워도 시간순으로 함께 읽으면 양립하지 않을 수 있다는 점이다.'' 수정 제안: ``하지만 이 방식은 같은 장면의 CoC들이 개별적으로는 자연스러워도 시간순으로 함께 읽으면 양립하지 않을 수 있다는 문제점이 있다.

\end{reviewercomment}

\begin{authorresponse}

제안하신 취지를 반영하여 해당 문단을 다시 구성했습니다. CoC가 행동의 이유는 말하지만 그 요구가 언제까지 유지되고 무엇으로 끝나는지는 대개 밝히지 않는다는 점을 먼저 설명하고, 앞선 정지 의무와 녹색 신호에 따른 진행을 함께 읽을 때 두 문장이 양립하지 않을 수 있음을 예로 보였습니다. 이 예는 자연 자료에서 발견한 사례가 아니라 설명을 위한 예시로 제시했습니다.

\end{authorresponse}

\begin{manuscriptrevision}{\mloc{sec:intro}; 발췌: \quotepages{R2-1.1}}
본문의 실제 수정 문장은 다음과 같습니다.

\begin{quote}
두 문장은 각각 자연스럽지만 함께 읽으면 문제가 생긴다. 보행자가 아직 횡단 중이라면 두 문장은 양립하지 않는다.
\end{quote}
\end{manuscriptrevision}

\end{reviewitem}

\clearpage

\begin{reviewitem}{R2-1.2. 불일치의 영향 문장}

\begin{reviewercomment}

기존 문장: ``이러한 불일치는 학습 자료의 행동 요구를 모호하게 하고 같은 자료를 다시 검사했을 때 일관된 결과를 얻기 어렵게 한다.'' 수정 제안: ``이러한 불일치 때문에 학습 자료의 행동 요구가 모호해지고, 같은 자료를 다시 검사했을 때 일관된 결과를 얻기 어렵다.

\end{reviewercomment}

\begin{authorresponse}

제안하신 표현을 따르되, 후반부는 같은 자료를 다시 검사할 때 어떤 조건과 근거를 적용해야 하는지 불분명해진다는 내용으로 구체화했습니다. 행동 요구의 모호성이 검토 동기임을 설명하되 이번 검사로 실제 학습 성능이 향상됐다고 주장하지 않습니다. 명시된 의무·증거의 일관성과 차량 안전 및 학습 효과를 구분했습니다.

\end{authorresponse}

\begin{manuscriptrevision}{\mloc{sec:intro}; 발췌: \quotepages{R2-1.2}}
본문의 실제 수정 문장은 다음과 같습니다.

\begin{quote}
이러한 불일치 때문에 학습 자료의 행동 요구가 모호해지고, 같은 자료를 다시 검사할 때 어떤 조건과 근거를 적용해야 하는지도 불분명해진다.
\end{quote}
\end{manuscriptrevision}

\end{reviewitem}

\clearpage

\begin{reviewitem}{R2-2. 논문 전반의 국문 표현}

\begin{reviewercomment}

논문 전체에 걸쳐 영문을 번역한 듯한 문장이 다수 있으므로, 자연스러운 국문 문장으로 수정할 필요가 있다.

\end{reviewercomment}

\begin{authorresponse}

서론뿐 아니라 관련 연구·자료·문제 정의·방법·평가·결과·논의·한계·결론을 새 근거에 맞춰 다시 교열했습니다. 주어와 판단 대상을 명시하고 ``확정 모순''ㆍ``근거 부족''ㆍ``양립''의 판정 범주와 ``조건 연결 모델''ㆍ``경계 확장 모델'', ``개발 후 동일 사례 재시험''의 용어를 일관되게 사용했습니다. 자동 문구 추출과 작성 증거, 주석 문제와 처리 오류를 한 문장에 섞지 않도록 정리했습니다.

독립 검증이나 합의 과정처럼 수행되지 않은 작업을 완료형으로 표현하지 않았습니다. 영문 초록과 결과 표의 분모도 국문 본문과 맞췄습니다.

\end{authorresponse}

\begin{manuscriptrevision}{\mloc{sec:intro}; 발췌: \quotepages{R2-2}}
본문의 실제 수정 문장은 다음과 같습니다.

\begin{quote}
본 연구는 이를 위해 연속 CoC를 검사하는 프레임워크를 제안한다(그림~1). 절차는 세 단계이다. 첫째, 같은 장면의 사건을 시간순으로 고정한다. 둘째, 각 사건에서 의무, 해제 조건, 조건별 증거와 후속 행동을 의미 프레임으로 구성하고, 각 항목이 원문의 어느 구절에서 왔는지 기록한다. 본 논문의 평가에서 이 프레임은 설계자가 직접 작성했다.
\end{quote}
\end{manuscriptrevision}

\end{reviewitem}

\clearpage

\begin{reviewitem}{R2-3. UPPAAL 명칭}

\begin{reviewercomment}

UPPAAL은 무엇의 약자인가?

\end{reviewercomment}

\begin{authorresponse}

UPPAAL은 Uppsala와 Aalborg의 이름을 연결한 도구명이며 임의의 일반 영어 구절로 풀어 쓰는 약어가 아닙니다. 개발기관의 공식 소개에서 두 대학의 공동 개발과 시간 오토마타 기반 도구라는 설명을 확인하고 최초 등장에 명칭과 기능을 덧붙였습니다. 관련 연구에서는 공식 질의 문서와 기존 튜토리얼을 인용했습니다. 도구가 시간 오토마타를 지원하는 사실과 본 중심 모델에 정량 시간 시계가 없는 사실도 구분했습니다.

\end{authorresponse}

\begin{manuscriptrevision}{\mloc{sec:intro}; 발췌: \quotepages{R2-3}}
본문의 실제 수정 문장은 다음과 같습니다.

\begin{quote}
UPPAAL은 웁살라(Uppsala) 대학과 올보르(Aalborg) 대학이 공동 개발한 시간 오토마타 기반 모델 검사 도구이며, 이름은 두 대학명을 결합한 것이다~[17, 18, 19].
\end{quote}
\end{manuscriptrevision}

\end{reviewitem}

\clearpage

\begin{reviewitem}{R2-4. 연구 질문과 기여의 연결}

\begin{reviewercomment}

서론의 연구 질문은 본 논문의 연구 동기(motivation)로 보인다. 그렇다면 연구 질문에 대한 답이 본 논문의 기여에 해당하는지를 명확하게 기술해야 한다.

\end{reviewercomment}

\begin{authorresponse}

연구 질문을 두 가지로 명시했습니다. RQ1은 의무와 해제 근거를 연결해 확정 모순과 근거 부족을 구분하고 최초 문제를 추적할 수 있는가, RQ2는 장면 경계에서 의무와 증거를 유지하는 정책이 판정을 어떻게 바꾸는가입니다. 서론의 네 기여에는 대응하는 연구 질문을 표시했습니다. 중간표현과 판정 규칙은 RQ1, 경계 입력 조건은 RQ2, UPPAAL 구현과 비교·경계·전수 시험은 두 질문 모두에 대응합니다. 결론도 두 질문의 답을 차례로 정리하도록 다시 썼습니다. 평가 결과 자체를 독립적인 알고리즘 기여로 세지 않습니다.

% Supplemental evidence 2026-10-12
현재 첫 질문은 의무·해제 근거 연결과 문제 추적으로, 144건 비교 및 동일 사례 재시험에 대응합니다. 둘째 경계 질문은 유지/종료·근거 만료·재활성화 정책으로, 별도 81건에 대응합니다. 유한 범위 전수 대조는 이 의미 명세의 구현 적합성을 보강하며 별도의 자연 성능 기여로 계산하지 않습니다.
% End supplemental evidence

\end{authorresponse}

\begin{manuscriptrevision}{\mloc{sec:intro}; 발췌: \quotepages{R2-4}}
본문의 실제 수정 문장은 다음과 같습니다.

\begin{quote}
본 논문의 연구 질문은 다음 두 가지이다. \textbf{RQ1:} 의무와 해제 근거를 연결하면 행동 전환에서 확정 모순과 근거 부족을 구분하고, 최초 문제 사건과 관련 의무를 추적할 수 있는가? \textbf{RQ2:} 장면 경계에서 대상 동일성, 적용 범위와 증거 유효성이 주어졌을 때, 의무와 증거를 유지하는 정책과 경계에서 초기화하는 정책은 판정을 어떻게 다르게 만드는가?
\end{quote}
\end{manuscriptrevision}

\end{reviewitem}

\clearpage

\begin{reviewitem}{R2-5. 서론 구조와 간결성}

\begin{reviewercomment}

서론에서 연구 질문에 대한 구체적인 해결 방법론이 구조적으로 설명되지 않은 것으로 보인다. 전체적으로 장황하므로 간결하고 명확하게 다시 기술할 필요가 있다.

\end{reviewercomment}

\begin{authorresponse}

전체 흐름 그림에는 자동 텍스트와 작성 프레임의 별도 입력 경로, 검사기와 문제 추적 출력을 표시했습니다.

예시에서 A 미해제 확인은 확정 모순, A 해제 미상은 근거 부족, A 해제 확인은 양립이라는 구분을 먼저 설명합니다. 상세 알고리즘·입력 사례·질의는 방법 절에 배치했습니다.

서론은 구조적 동기와 질문·기여를 유지하고, 단계별 상세 설명과 모델 도식은 방법 절에 보강했습니다. 반복되는 범위 제한은 한계 절에 모아 핵심 설명의 공간을 확보했습니다.

% Supplemental evidence 2026-10-12
서론을 배경·문제·실제 기록의 장면 경계 동기·기존 연구·접근 방법·연구 질문·기여·평가 범위·논문 구성의 순서로 갱신했습니다. 네 평가의 핵심 결과를 구분해 요약하고, 경계 정책·배열·질의와 상세 수치는 방법과 결과 절에서 설명했습니다.
% End supplemental evidence

\end{authorresponse}

\begin{manuscriptrevision}{\mloc{sec:intro}; 발췌: \quotepages{R2-5}}
본문의 실제 수정 문장은 다음과 같습니다.

\begin{quote}
본 연구는 이를 위해 연속 CoC를 검사하는 프레임워크를 제안한다(그림~1). 절차는 세 단계이다. 첫째, 같은 장면의 사건을 시간순으로 고정한다.
\end{quote}
\end{manuscriptrevision}

\end{reviewitem}

\clearpage

\begin{reviewitem}{R2-6. 관련 연구의 기술 비교}

\begin{reviewercomment}

2장 「관련 연구」에는 관련 연구 동향이나 기존 연구 결과의 문제점, 그리고 이를 보완할 내용이 제시되어야 하나, 현재는 충분히 서술되지 않았다. 이에 대한 보완이 필요하다.

\end{reviewercomment}

\begin{authorresponse}

관련 연구를 원문·공식 자료로 확인한 기술 비교로 보완했습니다. 설명 충실성 연구(Is VLA Reasoning Faithful?)와 VLADriveBench의 관측·개입 검증도 저장 주석의 조건 연결 검사와 비교했습니다. DriveLM/DriveVLM의 설명·계획 생성, runtime verification의 관측 경로 검사, LTL/TLTL의 유한 접두 경로 세 값 의미, UPPAAL의 상태 속성 검사와 Scenic/VerifAI의 상황 생성·분석을 입력·표현·출력으로 비교했습니다.

2장을 추론 기반 주행과 추론 충실성, 교통 규칙의 형식화와 검증, 자연어 요구의 형식화, 실행 중 검증과 모델 검사, 본 연구의 위치의 다섯 소절로 재구성했습니다. 사고 과정 설명의 충실성(Turpin 등, Lanham 등)과 VLA 추론–궤적 일관성 감시(Drive the Thoughts), 교통 규칙의 시간 논리 형식화와 규칙 검증(Esterle 등, Maierhofer 등, Wan 등), 자연어–시간 논리 변환(NL2TL, nl2spec)을 추가했습니다. 각 소절 끝에 해당 연구가 다루지 않는 부분, 즉 사람이 정한 규칙이나 개별 요구 문장이 아니라 자동 생성된 연속 주석에서 의무와 해제 근거를 연결하는 문제를 보완 방향으로 제시했습니다.

이전 상태 기억과 세 값·최초 위반이 선행 기술임을 명시했습니다. 기존 방법이 CoC 검사를 할 수 없다고 단정하지 않고, 본 논문이 원문 구절과 의무·조건·증거·문제 기록의 연결을 실증한다는 범위로 보완 방향을 설명했습니다. 문헌 비교를 직접 수행하지 않은 성능 우위로 바꾸지 않았습니다.

\end{authorresponse}

\begin{manuscriptrevision}{\mloc{sec:related-position}; 발췌: \quotepages{R2-6}}
본문의 실제 수정 문장은 다음과 같습니다.

\begin{quote}
기존 연구는 한 장면의 설명과 행동의 정렬, 사람이 정한 규칙의 형식화와 검증, 개별 요구 문장의 형식화, 그리고 형식 명세에 대한 경로 감시를 다룬다. 본 연구는 이들 사이에 남은 질문, 즉 자동 생성된 연속 주석에서 어떤 문장이 의무를 만들고 어떤 증거가 그것을 해제하는지를 연결하고 증거가 없는 경우를 모순과 구분하는 문제를 다룬다.
\end{quote}
\end{manuscriptrevision}

\end{reviewitem}

\clearpage

\begin{reviewitem}{R2-7. 기존 연구 대비 기여 범위}

\begin{reviewercomment}

관련 연구 및 기존 연구와 비교하여 본 논문 연구 결과의 독창성이 명확하게 제시되지 않은 것으로 보인다. 이에 대한 보완이 필요하다.

\end{reviewercomment}

\begin{authorresponse}

기여를 상태 기억이나 모델 검사 기법 자체의 발명으로 표현하지 않았습니다. 본 논문은 CoC 원문 구절에서 의무 ID·해제 조건·조건별 증거·행동 전환·최초 문제 기록을 연결해 검토 가능한 입력과 출력으로 만드는 데 초점을 둡니다. 보행자 세 변형의 끝까지 연결된 예시와 실제 질의/출력 대조로 구체화했습니다.

단순 FSM도 성공한 사실을 유지하여 UPPAAL 정확도 우위, 최초·최고 또는 기존 기술로 불가능하다는 주장은 하지 않았습니다. 현재 근거는 작성 시나리오에서의 방법 실증이며 일반적인 독창성·성능 주장은 제한했습니다.

\end{authorresponse}

\begin{manuscriptrevision}{\mloc{sec:intro}; 발췌: \quotepages{R2-7}}
본문의 실제 수정 문장은 다음과 같습니다.

\begin{quote}
원문 구절, 의무, 해제 조건과 조건별 증거를 연결하는 중간표현을 제시한다 (RQ1).
\end{quote}
\end{manuscriptrevision}

\end{reviewitem}

\section*{맺음말}

심사위원님의 의견에 따라 서론을 문제, 실제 기록의 동기, 기존 연구의 한계, 접근 방법, 연구 질문, 기여, 평가 범위의 순서로 다시 구성하고 연구 질문과 기여의 대응을 명시했습니다. 관련 연구를 다섯 연구군으로 재구성하여 기존 기법과 본 연구의 기여를 구분했으며, 원고 전반의 국문 표현과 판정 범주·모델 명칭을 정리했습니다.

\end{document}
